\documentclass[11pt,a4paper]{article}
\usepackage[margin=26mm]{geometry}
\usepackage{amsmath,amssymb,amsthm,mathtools,booktabs,microtype}
\usepackage[T1]{fontenc}
\usepackage{lmodern}
\usepackage{xurl}
\usepackage[hidelinks]{hyperref}
\newtheorem{theorem}{Theorem}
\newtheorem{proposition}[theorem]{Proposition}
\newtheorem{lemma}[theorem]{Lemma}
\newtheorem{corollary}[theorem]{Corollary}
\newcommand{\N}{\mathbb N}
\newcommand{\Z}{\mathbb Z}
\newcommand{\R}{\mathbb R}
\newcommand{\C}{\mathbb C}
\newcommand{\E}{\mathbb E}
\newcommand{\PP}{\mathbb P}
\newcommand{\eps}{\varepsilon}
\newcommand{\norm}[1]{\left\lVert#1\right\rVert}
\newcommand{\abs}[1]{\left|#1\right|}
\DeclareMathOperator{\Geom}{Geom}
\setlength{\parskip}{3pt}
\hypersetup{pdftitle={Fourier decay for truncated Syracuse affine laws},pdfauthor={Lucas Valbuena}}
\title{Fourier decay for truncated Syracuse affine laws}
\author{Lucas Valbuena}
\date{7 October 2026}

\begin{document}
\maketitle
\begin{abstract}
Let \(a_1,\ldots,a_k\) be independent positive geometric random variables
with \(\PP(a_i=a)=2^{-a}\), and set \(S_j=a_1+\cdots+a_j\).
We study the affine law
\(\Phi_k(z)=\sum_{j=1}^k3^{j-1}2^{-S_j}+3^k2^{-S_k}z\)
modulo \(3^n\), with \(k\le n\).
For every fixed \(A\in\N\), its primitive Fourier coefficients are bounded
by \(C_A((n-k+2)/n)^A\), uniformly in the seed and after multiplication
by any unit-modulus function of \(S_k\).
Consequently, conditioning on any total-exponent slice of probability at
least \(k^{-B}\) preserves decay faster than every fixed inverse power of
\(k\) throughout \(n-k\le k^\beta\), for each fixed \(0<\beta<1\).
The event \(S_k=2k\) satisfies this conclusion without an additional mass
hypothesis. The proof combines a fixed-dimensional Subspace Theorem
obstruction to finite black-triangle covers, a qualitative renewal
estimate above the entropy threshold, and a conductor recursion.
The constants are ineffective. The results concern finite random affine
laws and do not prove the Collatz conjecture.
\end{abstract}

\section{The random affine law and the main estimates}

Throughout, \(\N=\{0,1,2,\ldots\}\). The variables \(a_i\) are independent
and take positive integer values with \(\PP(a_i=a)=2^{-a}\).
Put \(S_0=0\) and \(S_j=\sum_{i=1}^ja_i\). Since two is a unit modulo
\(3^n\), the expression
\begin{equation}\label{eq:affine}
 \Phi_k(z)=\sum_{j=1}^k3^{j-1}2^{-S_j}+3^k2^{-S_k}z
 \quad\text{in }\Z/3^n\Z
\end{equation}
is defined for every \(n,k\in\N\) and \(z\in\Z/3^n\Z\).
This uses the reversed-variable convention for the Syracuse affine
composition. Reversing an iid exponent list leaves both its law and its
total exponent unchanged. We write
\[
 \chi_{n,\xi}(x)=\exp(-2\pi i\xi x/3^n).
\]
A frequency is \emph{primitive} if \(3\nmid\xi\).

\begin{theorem}[Conductor-deficit estimate]\label{thm:main}
For every \(A\in\N\) there is a constant \(C_A>0\) such that
\begin{equation}\label{eq:main}
 \abs{\E\bigl[\chi_{n,\xi}(\Phi_k(z))\tau(S_k)\bigr]}
 \le C_A\left(\frac{n-k+2}{n}\right)^A
\end{equation}
for all integers \(n\ge1\), \(0\le k\le n\), all primitive
\(\xi\in\Z/3^n\Z\), all \(z\in\Z/3^n\Z\), and every function
\(\tau:\N\to\C\) with \(\abs{\tau(s)}=1\) for all \(s\).
\end{theorem}

The constant precedes every parameter on the right of this statement.
In particular, the estimate is uniform when the deficit \(d=n-k\)
grows. It gives no power of \(n\) when \(d/n\) stays bounded away from
zero: the exponent \(A\) is fixed, and its constant is uncontrolled as
\(A\) varies.

Let \(p_k(s)=\PP(S_k=s)\). Conditioning below means the actual source
law restricted to \(S_k=s\) and divided by \(p_k(s)\).

\begin{theorem}[Conditioning on total exponent]\label{thm:conditional}
For each \(A\in\N\), one constant \(C_A>0\) works simultaneously for
all parameters in Theorem~\ref{thm:main} and every \(s\) with
\(p_k(s)>0\), giving
\begin{equation}\label{eq:conditional}
 \abs{\E[\chi_{n,\xi}(\Phi_k(z))\mid S_k=s]}
 \le \frac{C_A}{p_k(s)}\left(\frac{n-k+2}{n}\right)^A.
\end{equation}
Moreover, for every \(q,B\in\N\) and \(0<\beta<1\), there is
\(C_{q,B,\beta}>0\) such that
\begin{equation}\label{eq:massfloor}
 \abs{\E[\chi_{k+d,\xi}(\Phi_k(z))\mid S_k=s]}
 \le C_{q,B,\beta}k^{-q}
\end{equation}
whenever \(k\ge1\), \(d\in\N\), \(d\le k^\beta\), \(\xi\) is
primitive modulo \(3^{k+d}\), \(z\in\Z/3^{k+d}\Z\), and
\(p_k(s)\ge k^{-B}\).
\end{theorem}

\begin{corollary}[The central slice]\label{cor:central}
For every \(q\in\N\) and \(0<\beta<1\), there is \(C_{q,\beta}>0\)
such that
\[
 \abs{\E[\chi_{k+d,\xi}(\Phi_k(z))\mid S_k=2k]}
 \le C_{q,\beta}k^{-q}
\]
for every \(k\ge4\), \(d\le k^\beta\), every primitive frequency
modulo \(3^{k+d}\), and every seed \(z\) in that group.
\end{corollary}

The proof of the corollary uses the elementary, uniform lower bound
\(p_k(2k)\ge1/(4k)\). No local-limit estimate is an input.

\subsection{Relation to preceding work}

Tao's Proposition~1.17 gives the primitive Fourier estimate for the
full-length Syracuse offset law, and his Section~7 develops the paired
source and renewal geometry used here~\cite{Tao}.
At \(k=n\), the seed term vanishes modulo \(3^n\).
We do not claim this endpoint as a new result.
The distinction in \eqref{eq:main} is its dependence on a growing
conductor deficit. A bound with an exponential loss in \(d\) does not
give the range \(d\le k^\beta\).

Si's Theorem~H states uniform conditional decay at logarithmic
oversampling~\cite{Si}. Corollary~\ref{cor:central} supplies every fixed
polynomially sublinear oversampling range on the central slice;
Theorem~\ref{thm:conditional} covers all slices with the stated mass
floor. This comparison concerns those precise parameter ranges, and
does not assert a uniform theorem at a fixed positive oversampling
ratio. We do not use results from Si's preprint in the proof.

\section{Excluding a fixed number of black triangles}\label{sec:arithmetic}

Fix \(0\le\eps<1/4\) and a primitive frequency \(\xi\) modulo \(3^n\).
At a lattice point \((j,l)\in\Z_{>0}\times\Z\), let
\[
 \Theta_{n,\xi}(j,l)=\xi3^{2j-2}2^{1-l}\pmod{3^n}.
\]
The point is black if the absolute value of the signed representative
of \(\Theta_{n,\xi}(j,l)\), divided by \(3^n\), is at most \(\eps\).
Otherwise it is white. An integer-height triangle with top-left point
\((j_i,T_i)\) and height \(h_i\in\N\) consists of the lattice points
\begin{equation}\label{eq:triangle}
 j\ge j_i,\quad l\le T_i,\quad
 (j-j_i)\log9+(T_i-l)\log2\le h_i\log2.
\end{equation}
Calling a triangle black means that every lattice point in it is black.

Consider \(r+1\) such triangles, indexed by \(0\le i\le r\), with
\(2(j_i-1)<n\). Suppose there are \(m_i,u_i\in\N\) and
\(e_i,g_i\in\Z\) satisfying
\begin{align}\label{eq:schedules}
 j_{i+1}&=j_i+m_i+u_i,&
 T_{i+1}-T_i&=h_{i+1}+g_i &&(i<r),\\
 h_i&=4m_i+e_i,& \sum_{i=0}^r m_i&\le n.\nonumber
\end{align}
Define the error and span by
\begin{equation}\label{eq:error-span}
 E=\sum_{i=0}^r\abs{e_i}+\sum_{i=0}^{r-1}(u_i+\abs{g_i}),
 \qquad H=h_0+T_r-T_0.
\end{equation}
The vertical gaps are signed, and no ordering of the original tops is
assumed.

\begin{proposition}[Finite-cover obstruction]\label{prop:cover}
Fix \(R\in\N\), \(0\le\eps<1/4\), \(\eta>0\), and a function
\(w:\N\to\R\) such that, for every \(\gamma>0\),
\(w(n)\le\gamma n\) eventually. There is \(N\) such that every family
above with \(n\ge N\), \(r<R\), and \(E\le w(n)\) satisfies
\[
 H\log2<n\log3+\eta n.
\]
The cutoff is uniform over the primitive frequency and all triangle
locations. It need not be effective.
\end{proposition}

\begin{proof}
We first treat a putative sequence of counterexamples for which every
\(m_i\ge\delta n\), with fixed \(\delta>0\) and fixed \(r\).
Write \(q_i=3^{n-2j_i+2}\), and let \(a_i\) be the signed representative
of \(\xi2^{1-T_i}\) modulo \(q_i\). Primitivity gives \(a_i\ne0\).
Successive doublings down the left column cannot wrap while remaining
black when \(\eps<1/4\). Therefore
\begin{equation}\label{eq:black-numerator}
 \abs{a_i}\le\eps q_i2^{-h_i}.
\end{equation}
The height relation gives
\(D_i=T_{i+1}-T_i=h_{i+1}+g_i>0\) eventually.
The common frequency then gives exact integer carries \(z_i\) with
\begin{equation}\label{eq:carry}
 q_i=9^{m_i+u_i}q_{i+1},\qquad
 a_i=2^{D_i}a_{i+1}-z_iq_{i+1}.
\end{equation}
In particular,
\[
 \abs{z_i}\le\eps2^{g_i}
       +\eps9^{m_i+u_i}2^{-h_i}=\exp(o(n))
\]
as an upper bound, using \(e_i,g_i,u_i=o(n)\).

Put \(D=\sum_{i<r}D_i\), \(b=a_0\), and
\[
 X_0=2^Da_r,\qquad
 X_{i+1}=z_iq_{i+1}2^{D_0+\cdots+D_{i-1}}\quad(i<r).
\]
Unrolling \eqref{eq:carry} gives
\begin{equation}\label{eq:residual}
 b=X_0-\sum_{i<r}X_{i+1},\qquad
 0<\frac{\abs b}{2^D}\le\eps3^n2^{-H}.
\end{equation}
The single-triangle case, and the case of no nonzero carry, contradict
the entropy excess directly. Otherwise pass to a subsequence with a
fixed nonempty set of nonzero carry indices, and list the corresponding
coordinates as \(Y_1,\ldots,Y_t\), in their original order.

For the integer vector \(x=(X_0,Y_1,\ldots,Y_t)\), use the forms
\(X_0-\sum_jY_j,Y_1,\ldots,Y_t\) at infinity and the coordinate forms
at both primes two and three. Each system has determinant one.
Writing \(\abs{p}_p=p^{-1}\), their full local product is at most
\begin{equation}\label{eq:local-product}
 \left(\prod_{z_i\ne0}\abs{z_i}\right)\frac{\abs b}{2^D}
 \le \exp(-\eta n+o(n)).
\end{equation}
Indeed, the archimedean, two-adic and three-adic factors of a carry
coordinate together contribute at most \(\abs{z_i}\); the two-adic
factor of \(X_0\) is at most \(2^{-D}\).
All coordinates have absolute value at most \(\exp(Cn)\) for a fixed
\(C>0\). Taking \(\kappa=\eta/(4C)\), the product in
\eqref{eq:local-product} is eventually smaller than
\(\norm{x}_\infty^{-\kappa}\). The rational Subspace
Theorem~\cite{Schlickewei} thus places
these vectors in finitely many proper rational subspaces.

None of those subspaces can contain an infinite subsequence.
The exact schedules and \(16>9\) give
\begin{equation}\label{eq:ratios}
 \frac{b}{Y_j}\longrightarrow0,
 \qquad \frac{Y_i}{Y_j}\longrightarrow0\quad(i<j).
\end{equation}
For example, the first ratio is bounded in absolute value by
\((9/16)^{m_0+\cdots+m_i}\exp(o(n))\) for the carry indexed by \(i\).
The ratio of successive nonzero carry coordinates has the same form
with the intervening block sum. Their main exponents are linear in
\(n\), whereas all carry and scheduling losses are subexponential.
If a fixed rational relation
\(c_0X_0+\sum_jc_jY_j=0\) held along an infinite subsequence, then
\eqref{eq:residual} would give
\[
 c_0b+\sum_j(c_0+c_j)Y_j=0.
\]
Dividing by the last coordinate with nonzero coefficient and using
\eqref{eq:ratios} shows that every \(c_0+c_j\) vanishes. Since
\(b\ne0\), all coefficients then vanish, a contradiction.

For a general counterexample sequence, fix the number of triangles
on a subsequence and make every \(m_i/n\) converge. Delete the indices
with zero limit. At least one remains, since otherwise \(H=o(n)\).
The deleted lengths and heights total \(o(n)\). Joining the schedules
across each deleted segment adds these lengths to the horizontal
omissions and these heights to the signed vertical gaps. The total
error remains \(o(n)\), and the original and retained spans differ
in absolute value by \(o(n)\). The retained original
triangles have \(m_i\ge\delta n\) for one \(\delta>0\), and retain a
positive entropy excess. This reduces to the contradiction just proved.
Finally, failure of a uniform cutoff would supply an unbounded
counterexample sequence with the same fixed \(R,\eps,\eta,w\).
\end{proof}

\section{Qualitative decay above the entropy threshold}\label{sec:qualitative}

Pair the exponents and write
\[
 b_j=a_{2j-1}+a_{2j},\qquad L_j=b_1+\cdots+b_j.
\]
The \(b_j\) are independent with
\(\PP(b_j=b)=(b-1)2^{-b}\) for \(b\ge2\). For fixed pair sums,
the two splits of \(b_j=3\) are equiprobable. Their contributions to
the offset differ by
\(2\cdot3^{2j-2}2^{-L_j}\). Define
\[
 f_{n,\xi}(j,l)=
 \left|\frac{\chi_{n,\xi}(5\cdot3^{2j-2}2^{-l})+
                   \chi_{n,\xi}(7\cdot3^{2j-2}2^{-l})}{2}\right|
\]
and the nonnegative envelope
\begin{equation}\label{eq:envelope}
 Q_{n,\xi}(J)=\E\prod_{j=1}^J
 \begin{cases}f_{n,\xi}(j,L_j),&b_j=3,\\1,&b_j\ne3.\end{cases}
 \qquad
 Q(n,J)=\sup_{3\nmid\xi}Q_{n,\xi}(J).
\end{equation}
Thus \(0\le Q(n,J)\le1\). If \(J=\lfloor k/2\rfloor\),
conditioning on the pair sums and, for odd \(k\), the last exponent,
gives
\begin{equation}\label{eq:pair-majorant}
 \abs{\E[\chi_{n,\xi}(\Phi_k(z))\tau(S_k)]}
 \le Q_{n,\xi}(J).
\end{equation}
The total exponent, seed term, and terminal phase are constant during
this conditional averaging. The factors from other pair sums and the
unpaired exponent have absolute value at most one.

Fix a sufficiently small absolute \(\eps>0\), as in the renewal
estimate below. Let \(W_J\) count the indices \(j\le J\) with
\(b_j=3\) for which \((j,L_j)\) is white. The cosine formula and
the definition of white give
\begin{equation}\label{eq:laplace}
 Q_{n,\xi}(J)\le\E e^{-\eps^3W_J}
 \le\PP(W_J\le T)+e^{-\eps^3T}\qquad(T\in\N).
\end{equation}

\subsection{The finite renewal estimate}

The points \((j,L_j)\) for which \(b_j=3\), in their original order,
form a renewal path. Its increments are the independent blocks ending
at the next occurrence of \(b=3\). Each horizontal increment is
positive, each vertical increment is nonnegative, and their mean
slope is four. We use the black-triangle decomposition and renewal
estimate of Tao's Section~7~\cite{Tao}, in the following finite form.

Set \(C=\lfloor n/2\rfloor\), and use a list of \(C+1\) renewal
increments. Its first increment is the starting point of a path with
\(C\) subsequent steps. This list contains every raw success with
horizontal coordinate at most \(C\). The canonical black triangles
cover the black points in \(1\le j\le C\). Starting with the first
black point of this path, record its triangle; the next recorded point
is the first subsequent black point above the preceding triangle's
top. Continue with the same rule at path times less than \(C\).
Stops are numbered from one, so
the first stop can occur at path time zero. A missing stop is distinct
from a stop at zero.

Write \(t_R\) for the path time of stop \(R\), if it exists, and
let \(V(t)\) count the white points at positive path times through
\(t\), restricted to \(j\le C\). The finite stopped estimate is
\begin{equation}\label{eq:moment}
 \E\left[\mathbf1_{\{t_R\text{ exists}\}}
                e^{-V(t_R)+\eps R}\right]\le e^\eps.
\end{equation}
It holds for every \(R\ge1\), uniformly over the initial renewal
point and the primitive frequency. Averaging over the first increment
therefore preserves it.

For precision, \eqref{eq:moment} uses the full cutoff \(C\), even
when the desired horizontal horizon \(J\) is smaller. It is the
stopped, zero-on-absence version of the renewal argument underlying
Tao's Lemma~7.9. The one-step argument uses the localized
first-passage mass from Lemma~7.7 and the separation of distinct black
triangles. Choose the localization radius first, then choose \(\eps\)
so that its separation scale contains the resulting collar. The
one-step estimate applies uniformly at the restart point. We give the
finite induction explicitly to specify its cutoff convention.

Let \(G_R\) denote the random variable in \eqref{eq:moment}. Partition
according to the first-entry event \(A_{p,x}\), with first stop at
path time \(p\), point \(x\), and canonical triangle \(\Delta\).
Let \(K\ge1\) be the first additional time above the top of
\(\Delta\), and \(U\) its endpoint. Write \(w(U)=e^{-1}\) when
\(U\) is white and lies in \(j\le C\), and \(w(U)=1\) otherwise.
For \(R\ge2\), the stopping rule and white counts give
\begin{equation}\label{eq:one-step}
 G_R\le e^{-\max(p-1,0)+\eps}\,w(U)\,
              G_{R-1}(U;\text{fresh tail})
 \quad\text{on }A_{p,x}.
\end{equation}
If \(p+K\ge C\), the left side is zero. On the other branch, there
are at least \(\max(p-1,0)\) positive-time white points before the
first entry, and the first-exit point supplies the factor \(w(U)\).
The restarted variable uses a fresh list of \(C\) increments with
the same absolute horizontal cutoff \(C\). It does not replace that
cutoff by \(C-p-K\). Appending independent increments is legitimate
because the original event and statistic depend on the original
finite prefix only; the first-passage split then leaves an independent
fresh tail.

Here is the geometric input for the exit weight. For an entry at depth
\(s\) below its triangle's top, the relative first-passage endpoint
\((r,l)\) has probability at least \(1/2\) in the window
\[
 \abs{r-s/4}<L\sqrt{1+s},\qquad s<l<s+L,
\]
for one absolute \(L\), by the localized first-passage estimate.
Put \(a=\log2/\log9-1/4>0\) and
\(K_L=1+a+L^2/(4a)\). This window lies within distance
\(\sqrt{K_L^2+L^2}\) of a top-row point of the starting triangle.
The canonical geometry gives separation at least
\(\rho=\frac1{10}\log(1/\eps)\) between distinct triangles and a
right boundary at most \(n/2-\rho\) for every triangle.
Choose \(\eps\) small enough that
\(K_L^2+L^2\le\rho^2\) and \(\eps\le e^{-10}\).
Then the localized exit is outside its old triangle, inside the
horizontal cutoff, and in no other black triangle. It is white.
Consequently
\[
 \E w(U)\le c:=\frac{1+e^{-1}}2\le e^{-\eps}.
\]
Assuming \eqref{eq:moment} at \(R-1\), uniformly in the restart
point, conditional expectation in \eqref{eq:one-step} yields
\[
 \E[\mathbf1_{A_{p,x}}G_R]
 \le\PP(A_{p,x})e^{-\max(p-1,0)+2\eps}c
 \le\PP(A_{p,x})e^\eps.
\]
Sum over the countable first-entry events. Absence of a first entry
contributes zero; at \(R=1\) the pointwise bound is
\(G_1\le e^\eps\). This proves \eqref{eq:moment}, including an
entry at time zero.

Let \(E_R(J)\) be the event that stop \(R\) exists and its horizontal
coordinate is at most \(J\). On \(E_R(J)\), every white point counted
by \(V(t_R)\) is among those counted by \(W_J\). Markov's inequality
in \eqref{eq:moment} consequently gives
\begin{equation}\label{eq:stopped-tail}
 \PP\bigl(E_R(J)\cap\{W_J\le T\}\bigr)
 \le e^{T+\eps-\eps R}.
\end{equation}
This deduction requires no uniform estimate for a process restarted
with only \(J\) units of horizontal room.

\subsection{Good paths cannot use only finitely many triangles}

Take
\[
 H_n=\lceil\sqrt n\rceil,\qquad V_n=H_n,\qquad D_n=n^{3/4}.
\]
Call the renewal path good if its starting horizontal coordinate is
at most \(H_n\), its horizontal and vertical steps are at most
\(H_n,V_n\), and
\(\abs{l-4j}\le D_n\) at every path point with \(j\le J\).
Exponential-moment bounds for the independent blocks and for the raw
geometric sums give, for \(1\le2J\le n\),
\begin{equation}\label{eq:bad-path}
 \PP(\text{path is not good})\le B(n):=
 15(n+1)e^{-\log(21/20)\sqrt n}
       +2n e^{-\sqrt n/32}=o(1).
\end{equation}
These constants follow from the increment bound
\[
 \PP(\Delta j>H\text{ or }\Delta l>V)
 \le15e^{-\log(21/20)\min(H,V)}
\]
and the union bound over the first \(2J\) raw exponent sums, with
individual tail bound
\[
 2\exp\left[-\min\left(\frac{D^2}{32(2J)},\frac D8\right)\right].
\]

Fix \(R,T\), and assume the path is good, \(W_J\le T\), and
\(E_R(J)\) fails. The recorded black triangles before the first
crossing of \(J\) then form a family with at most \(R\) members.
There is a member once \(n\) is large: before encountering the first
black point, at most \(T\) white renewal points can advance the path
by at most \((T+1)H_n=o(n)\).

Write \((x_i,y_i)\), \(0\le i\le r\), for the recorded entry points,
and \(A_i\) for the corresponding original triangle tops. Re-anchor
each triangle at its entry, using the heights
\[
 h_i=A_i-y_i\quad(i<r),\qquad
 h_r=\min\{A_r-y_r,\,4(J-x_r)\},\qquad T_i=y_i+h_i.
\]
The new triangle with left edge \(x_i\), top \(T_i\), and height
\(h_i\) is contained in the original black triangle: add its defining
inequality to the inequality expressing membership of the entry
point in the original triangle. Set
\begin{align*}
 m_i&=x_{i+1}-x_i\quad(i<r),& m_r&=J-x_r,&u_i&=0,\\
 e_i&=h_i-4m_i,&g_i&=y_{i+1}-T_i\quad(i<r).
\end{align*}
These choices satisfy \eqref{eq:schedules} exactly.
Only points before the crossing were assumed to lie in the tube.
If \((j_-,l_-)\), \((j_+,l_+)\) straddle the crossing, use the
terminal height \(S=\min(4J,l_+)\). Positivity of vertical steps and
\(j_+-j_-\le H_n\) put this height within \(D_n+4H_n\) of \(4J\).
Put \(D'=D_n+4H_n\) to account for this endpoint.

The white gaps use at most \(T\) steps, with one exit overshoot for
each intervening triangle. Thus
\[
 \sum_{i<r}\abs{g_i}\le(T+r)V_n,\qquad
 x_0\le(T+1)H_n,\qquad A_r\ge S-(T+1)V_n.
\]
The tube identities and the chosen heights give
\(\abs{e_i}\le2D'+\abs{g_i}\) for \(i<r\), and
\(\abs{e_r}\le2D'+(T+1)V_n\). Consequently
\begin{align}\label{eq:path-cover}
 E&\le2(r+1)D'+(T+1)V_n+2(T+r)V_n,\\
 H&\ge4J-4(T+1)H_n-2D'-(T+1)V_n.\nonumber
\end{align}
Here \(H\) is the span from \eqref{eq:error-span}, distinct from
the step threshold \(H_n\). The vertical gaps remain signed, and
the retained horizontal lengths sum to \(J-x_0\le J\le n\).
Since \(r<R\) and \(R,T\) are fixed, the error and span loss in
\eqref{eq:path-cover} are \(o(n)\).

Suppose now that
\begin{equation}\label{eq:entropy}
 \log3<2\theta\log2,\qquad \theta n\le2J\le n.
\end{equation}
Then \eqref{eq:path-cover} implies
\(H\log2\ge2\theta n\log2-o(n)\).
Proposition~\ref{prop:cover}, with any fixed positive \(\eta\) smaller
than \(2\theta\log2-\log3\), excludes this family for all sufficiently
large \(n\). The cutoff is uniform in \(J\) and \(\xi\).
Thus a good path with \(W_J\le T\) must satisfy \(E_R(J)\), and
\eqref{eq:bad-path} and \eqref{eq:stopped-tail} yield
\begin{equation}\label{eq:few-white}
 \PP(W_J\le T)\le B(n)+e^{T+\eps-\eps R}
\end{equation}
for all sufficiently large \(n\), for each fixed \(R,T\).

\begin{proposition}[Uniform qualitative contraction]\label{prop:qualitative}
If \(\log3<2\theta\log2\), then, for every \(\delta>0\), there is
\(N\) such that
\[
 Q(n,J)\le\delta\qquad
 (n\ge N,\ \theta n\le2J\le n).
\]
\end{proposition}
\begin{proof}
Choose \(T\) with \(e^{-\eps^3T}<\delta/3\), then \(R\) with
\(e^{T+\eps-\eps R}<\delta/3\). Finally choose \(n\) beyond the
uniform cutoff in \eqref{eq:few-white} and with \(B(n)<\delta/3\).
Apply \eqref{eq:laplace}. No quantitative dependence of this cutoff
on \(R\) is needed.
\end{proof}

\section{The conductor recursion}\label{sec:recursion}

Split the product in \eqref{eq:envelope} after \(u\) pairs.
Conditional on the first \(u\) pair sums, its suffix has conductor
\(n-2u\), length \(J-u\), and primitive frequency multiplied by
the unit \(2^{-L_u}\). Independence and the supremum over primitive
frequencies imply
\begin{equation}\label{eq:recursion}
 Q(n,J)\le Q(n,u)Q(n-2u,J-u)
 \qquad(0\le u\le J,\ 2u<n).
\end{equation}

\begin{lemma}[Scalar bootstrap]\label{lem:bootstrap}
Suppose \(0\le Q(n,J)\le1\) on \(n\ge1\), \(2J\le n\), and
\eqref{eq:recursion} holds. If Proposition~\ref{prop:qualitative}
holds for some \(0<\theta<1\), then for every \(A\in\N\) there is
\(C_A>0\) such that
\begin{equation}\label{eq:paired-deficit}
 Q(n,J)\le C_A\left(\frac{n-2J+1}{n}\right)^A.
\end{equation}
\end{lemma}
\begin{proof}
Put \(\lambda=(1-\theta)/2\), and obtain the qualitative contraction
with \(\delta=\lambda^A\). For large \(n\), set
\(u=\lfloor(1-\lambda)n/2\rfloor\) and \(n'=n-2u\).
Then
\[
 1\le u,\qquad \theta n\le2u<n,\qquad
 \lambda n\le n'<\lambda n+2.
\]
Let \(D=n-2J+1\). If \(D\le\lambda n\), then \(u\le J\), and
the deficit is preserved by the recursion:
\(n'-2(J-u)+1=D\).
Strong induction on \(n\) gives
\[
 Q(n,J)\le\lambda^A C_A(D/n')^A
             \le C_A(D/n)^A.
\]
If \(D>\lambda n\), the trivial bound one suffices after taking
\(C_A\ge\lambda^{-A}\). The finitely many smaller conductors are
also absorbed in \(C_A\), since \(D\ge1\).
\end{proof}

\begin{proof}[Proof of Theorem~\ref{thm:main}]
Choose \(\theta=5/6\). The entropy inequality follows from
\(3^3=27<32=2^5\), and \(0<\theta<1\).
Apply Proposition~\ref{prop:qualitative} and Lemma~\ref{lem:bootstrap}.
For \(J=\lfloor k/2\rfloor\), one has
\(n-2J+1\le n-k+2\). The conclusion follows from
\eqref{eq:pair-majorant} and \eqref{eq:paired-deficit}.
\end{proof}

\section{Conditioned affine laws}\label{sec:conditioning}

\begin{proof}[Proof of Theorem~\ref{thm:conditional}]
For fixed \(s\), let \(\sigma_s(t)=1\) if \(t=s\) and \(-1\)
otherwise. Both \(1\) and \(\sigma_s\) are permitted terminal
phases, and
\[
 \E[\chi_{n,\xi}(\Phi_k(z))\mathbf1_{\{S_k=s\}}]
 =\tfrac12\E[\chi_{n,\xi}(\Phi_k(z))]
  +\tfrac12\E[\chi_{n,\xi}(\Phi_k(z))\sigma_s(S_k)].
\]
The triangle inequality and Theorem~\ref{thm:main}, followed by division
by \(p_k(s)>0\), prove \eqref{eq:conditional} with the same constant.

For \eqref{eq:massfloor}, choose an integer \(M\) with
\((1-\beta)M\ge q+B\). When \(n=k+d\), \(k\ge1\), and
\(d\le k^\beta\),
\[
 \frac{n-k+2}{n}=\frac{d+2}{k+d}
 \le 3k^{\beta-1}.
\]
Use \eqref{eq:conditional} with exponent \(M\) and
\(p_k(s)^{-1}\le k^B\). The resulting upper bound is
\(C_M3^M k^{(\beta-1)M+B}\le C_M3^Mk^{-q}\).
\end{proof}

\begin{proof}[Proof of Corollary~\ref{cor:central}]
Every positive composition of \(s\) into \(k\) parts has probability
\(2^{-s}\), and there are \(\binom{s-1}{k-1}\) such compositions.
Thus, for \(k\ge1\), \(s\ge1\),
\[
 p_k(s)=\binom{s-1}{k-1}2^{-s}.
\]
The usual convention makes this zero for \(s<k\).
The elementary central-binomial bound
\(4^k\le2k\binom{2k}{k}\), together with
\(\binom{2k}{k}=2\binom{2k-1}{k-1}\), gives
\[
 p_k(2k)=\binom{2k-1}{k-1}2^{-2k}\ge\frac1{4k}.
\]
For \(k\ge4\), this is at least \(k^{-2}\). Apply
Theorem~\ref{thm:conditional} with \(B=2\).
\end{proof}

\section{Formal proof and scope}\label{sec:formal}

The accompanying Lean sources prove the actual geometric-list law,
its pair-sum envelope, the canonical renewal estimate, the arithmetic
obstruction, and the conductor bootstrap. The conditioned law is
defined by filtering the source probability mass function on the
exact exponent sum, then mapping the affine offset; its Fourier
coefficient is proved equal to the normalized slice expectation.
It is not an abstract distribution assumed to have that identity.

The main theorem names, in namespace \texttt{Erdos1135.Tao}, are
\begin{itemize}
\item \path{trap_unconditional_affineTwistedExpectation_deficit_all_powers},
for the joint seed and terminal-phase estimate;
\item \path{trap_unconditional_seededAffinePMF_deficit_all_powers},
for the seeded affine probability mass function;
\item \path{trap_unconditional_conditionedAffinePMF_polynomial_deficit_mass_floor},
for Theorem~\ref{thm:conditional}'s polynomial deficit range; and
\item \path{trap_unconditional_centralConditionedAffinePMF_polynomial_deficit},
for Corollary~\ref{cor:central}.
\end{itemize}
The finite-cover argument uses the rational Subspace Theorem from
Ralf Stephan's \emph{Subspace-Theorems} formalization~\cite{SubspaceLean}.
The package
pins its source revision and the Lean and Mathlib versions, includes
the source dependency closure, and checks theorem dependencies for
unproved assumptions. The Subspace Theorem is proved in that closure;
it is not introduced as a new axiom.

Automated assistants were used for proof development, source review,
and preparation of this manuscript. The accompanying verification
records describe those checks; they do not constitute external peer
review.

These estimates concern a finite random affine model. They do not
bound the exceptional orbits of the deterministic Syracuse map.
In particular, Fourier decay on primitive frequencies does not by
itself provide the uniform integrability or integer-occupation upper
bound needed to transfer these estimates to every starting integer.
Neither convergence of every Collatz orbit nor absence of nontrivial
cycles is asserted here.

\begin{thebibliography}{9}
\bibitem{Tao}
T. Tao, \emph{Almost all orbits of the Collatz map attain almost
bounded values}, Forum of Mathematics, Pi \textbf{10} (2022), e12.
\url{https://arxiv.org/abs/1909.03562}.

\bibitem{Si}
Y. Si, \emph{A microcanonical phase transition for the Collatz affine
random model}, preprint,
4 May 2026, Theorem~H.
\url{https://zenodo.org/records/20027097}.

\bibitem{Schlickewei}
H. P. Schlickewei, \emph{The quantitative Subspace Theorem for number
fields}, Compositio Mathematica \textbf{82} (1992), 245--273.
\url{https://www.numdam.org/item/CM_1992__82_3_245_0/}.

\bibitem{SubspaceLean}
R. Stephan, \emph{Subspace-Theorems}, Lean formalization, revision
\path{d7a11cb1dcab3883fd1008db27a8382c95455b0a}.
\url{https://github.com/rwst/Subspace-Theorems}.
\end{thebibliography}
\end{document}
